\usepackage[cm]{fullpage}
\usepackage{color}
\usepackage{hyperref}

\hypersetup{breaklinks=true,%
pagecolor=white,%
colorlinks=true,%
linkcolor=cyan,%
urlcolor=MyDarkBlue}

\definecolor{MyDarkBlue}{rgb}{0,0.0,0.45}

%%%%%%%%%%%%%%%%%%%%%%%%%%
% Formatting parameters  %
%%%%%%%%%%%%%%%%%%%%%%%%%%

\newlength{\tabin}
\setlength{\tabin}{1em}
\newlength{\secsep}
\setlength{\secsep}{0.1cm}

\newlength{\projectspacing}
\setlength{\projectspacing}{1em}

\setlength{\parindent}{0in}
\setlength{\parskip}{0in}
    \setlength{\itemsep}{0in}
    \setlength{\topsep}{0in}
\setlength{\tabcolsep}{0in}

\definecolor{contactgray}{gray}{0.3}
\pagestyle{empty}

%%%%%%%%%%%%%%%%%%%%%%%%%%
%   Per-variant knobs    %
%%%%%%%%%%%%%%%%%%%%%%%%%%
% Overridden by a CV variant after \usepackage[cm]{fullpage}
\usepackage{color}
\usepackage{hyperref}

\hypersetup{breaklinks=true,%
pagecolor=white,%
colorlinks=true,%
linkcolor=cyan,%
urlcolor=MyDarkBlue}

\definecolor{MyDarkBlue}{rgb}{0,0.0,0.45}

%%%%%%%%%%%%%%%%%%%%%%%%%%
% Formatting parameters  %
%%%%%%%%%%%%%%%%%%%%%%%%%%

\newlength{\tabin}
\setlength{\tabin}{1em}
\newlength{\secsep}
\setlength{\secsep}{0.1cm}

\newlength{\projectspacing}
\setlength{\projectspacing}{1em}

\setlength{\parindent}{0in}
\setlength{\parskip}{0in}
    \setlength{\itemsep}{0in}
    \setlength{\topsep}{0in}
\setlength{\tabcolsep}{0in}

\definecolor{contactgray}{gray}{0.3}
\pagestyle{empty}

%%%%%%%%%%%%%%%%%%%%%%%%%%
%   Per-variant knobs    %
%%%%%%%%%%%%%%%%%%%%%%%%%%
% Overridden by a CV variant after \usepackage[cm]{fullpage}
\usepackage{color}
\usepackage{hyperref}

\hypersetup{breaklinks=true,%
pagecolor=white,%
colorlinks=true,%
linkcolor=cyan,%
urlcolor=MyDarkBlue}

\definecolor{MyDarkBlue}{rgb}{0,0.0,0.45}

%%%%%%%%%%%%%%%%%%%%%%%%%%
% Formatting parameters  %
%%%%%%%%%%%%%%%%%%%%%%%%%%

\newlength{\tabin}
\setlength{\tabin}{1em}
\newlength{\secsep}
\setlength{\secsep}{0.1cm}

\newlength{\projectspacing}
\setlength{\projectspacing}{1em}

\setlength{\parindent}{0in}
\setlength{\parskip}{0in}
    \setlength{\itemsep}{0in}
    \setlength{\topsep}{0in}
\setlength{\tabcolsep}{0in}

\definecolor{contactgray}{gray}{0.3}
\pagestyle{empty}

%%%%%%%%%%%%%%%%%%%%%%%%%%
%   Per-variant knobs    %
%%%%%%%%%%%%%%%%%%%%%%%%%%
% Overridden by a CV variant after \usepackage[cm]{fullpage}
\usepackage{color}
\usepackage{hyperref}

\hypersetup{breaklinks=true,%
pagecolor=white,%
colorlinks=true,%
linkcolor=cyan,%
urlcolor=MyDarkBlue}

\definecolor{MyDarkBlue}{rgb}{0,0.0,0.45}

%%%%%%%%%%%%%%%%%%%%%%%%%%
% Formatting parameters  %
%%%%%%%%%%%%%%%%%%%%%%%%%%

\newlength{\tabin}
\setlength{\tabin}{1em}
\newlength{\secsep}
\setlength{\secsep}{0.1cm}

\newlength{\projectspacing}
\setlength{\projectspacing}{1em}

\setlength{\parindent}{0in}
\setlength{\parskip}{0in}
    \setlength{\itemsep}{0in}
    \setlength{\topsep}{0in}
\setlength{\tabcolsep}{0in}

\definecolor{contactgray}{gray}{0.3}
\pagestyle{empty}

%%%%%%%%%%%%%%%%%%%%%%%%%%
%   Per-variant knobs    %
%%%%%%%%%%%%%%%%%%%%%%%%%%
% Overridden by a CV variant after \input{preamble.tex}, before the document.

% Whether a section is kept on a single page where possible. cv-ios.tex takes
% the default; cv-backend.tex sets this empty, matching its Overleaf original.
\newcommand{\sectionfilbreak}{\filbreak}

% \interests is deliberately NOT defined here. Each variant must define it, so
% a missing definition is a build error rather than a placeholder that could
% slip into a PDF you send out. See cv-ios.tex / cv-backend.tex.

%%%%%%%%%%%%%%%%%%%%%%%%%%
%  Template Definitions  %
%%%%%%%%%%%%%%%%%%%%%%%%%%

\newcommand{\lineunder}{\vspace*{-8pt} \\ \hspace*{-6pt} \hrulefill \\ \vspace*{-15pt}}
\newcommand{\name}[1]{\begin{center}\textsc{\Huge#1}\\\end{center}}
\newcommand{\program}[1]{\begin{center}\textsc{#1}\end{center}}
\newcommand{\contact}[1]{\begin{center}\color{contactgray}{\small#1}\end{center}}

\newenvironment{tabbedsection}[1]{
  \begin{list}{}{
      \setlength{\itemsep}{0pt}
      \setlength{\labelsep}{0pt}
      \setlength{\labelwidth}{0pt}
      \setlength{\leftmargin}{\tabin}
      \setlength{\rightmargin}{\tabin}
      \setlength{\listparindent}{0pt}
      \setlength{\parsep}{0pt}
      \setlength{\parskip}{0pt}
      \setlength{\partopsep}{0pt}
      \setlength{\topsep}{#1}
    }
  \item[]
}{\end{list}}

\newenvironment{nospacetabbing}{
    \begin{tabbing}
}{\end{tabbing}\vspace{-1.2em}}

\newenvironment{resume_header}{}{\vspace{0pt}}


\newenvironment{resume_section}[1]{
  \sectionfilbreak
  \vspace{2\secsep}
  \textsc{\large#1}
  \lineunder
  \begin{tabbedsection}{\secsep}
}{\end{tabbedsection}}

\newenvironment{resume_subsection}[2][]{
  \textbf{#2} \hfill {\footnotesize #1} \hspace{2em}
  \begin{tabbedsection}{0.5\secsep}
}{\end{tabbedsection}}

\newenvironment{subitems}{
  \renewcommand{\labelitemi}{-}
  \begin{itemize}
      \setlength{\labelsep}{1em}
}{\end{itemize}}

\newenvironment{resume_employer}[4]{
  \vspace{\secsep}
  \textbf{#1} \\ 
  \indent {\small #2} \hfill {\footnotesize#3 (#4)}
  \begin{tabbedsection}{0pt}
  \begin{subitems}
}{\end{subitems}\end{tabbedsection}}

%%%%%%%%%%%%%%%%%%%%%%%%%%
%    Contact details     %
%%%%%%%%%%%%%%%%%%%%%%%%%%
% Real phone/email live in contact.tex, which is gitignored so they are not
% published. A clone of the public repo has only contact.example.tex, hence
% the fallback -- without it a fresh clone would fail to build.
% Defines: \contactdetails

\IfFileExists{contact.tex}
  {\input{contact.tex}}
  {\input{contact.example.tex}}
, before the document.

% Whether a section is kept on a single page where possible. cv-ios.tex takes
% the default; cv-backend.tex sets this empty, matching its Overleaf original.
\newcommand{\sectionfilbreak}{\filbreak}

% \interests is deliberately NOT defined here. Each variant must define it, so
% a missing definition is a build error rather than a placeholder that could
% slip into a PDF you send out. See cv-ios.tex / cv-backend.tex.

%%%%%%%%%%%%%%%%%%%%%%%%%%
%  Template Definitions  %
%%%%%%%%%%%%%%%%%%%%%%%%%%

\newcommand{\lineunder}{\vspace*{-8pt} \\ \hspace*{-6pt} \hrulefill \\ \vspace*{-15pt}}
\newcommand{\name}[1]{\begin{center}\textsc{\Huge#1}\\\end{center}}
\newcommand{\program}[1]{\begin{center}\textsc{#1}\end{center}}
\newcommand{\contact}[1]{\begin{center}\color{contactgray}{\small#1}\end{center}}

\newenvironment{tabbedsection}[1]{
  \begin{list}{}{
      \setlength{\itemsep}{0pt}
      \setlength{\labelsep}{0pt}
      \setlength{\labelwidth}{0pt}
      \setlength{\leftmargin}{\tabin}
      \setlength{\rightmargin}{\tabin}
      \setlength{\listparindent}{0pt}
      \setlength{\parsep}{0pt}
      \setlength{\parskip}{0pt}
      \setlength{\partopsep}{0pt}
      \setlength{\topsep}{#1}
    }
  \item[]
}{\end{list}}

\newenvironment{nospacetabbing}{
    \begin{tabbing}
}{\end{tabbing}\vspace{-1.2em}}

\newenvironment{resume_header}{}{\vspace{0pt}}


\newenvironment{resume_section}[1]{
  \sectionfilbreak
  \vspace{2\secsep}
  \textsc{\large#1}
  \lineunder
  \begin{tabbedsection}{\secsep}
}{\end{tabbedsection}}

\newenvironment{resume_subsection}[2][]{
  \textbf{#2} \hfill {\footnotesize #1} \hspace{2em}
  \begin{tabbedsection}{0.5\secsep}
}{\end{tabbedsection}}

\newenvironment{subitems}{
  \renewcommand{\labelitemi}{-}
  \begin{itemize}
      \setlength{\labelsep}{1em}
}{\end{itemize}}

\newenvironment{resume_employer}[4]{
  \vspace{\secsep}
  \textbf{#1} \\ 
  \indent {\small #2} \hfill {\footnotesize#3 (#4)}
  \begin{tabbedsection}{0pt}
  \begin{subitems}
}{\end{subitems}\end{tabbedsection}}

%%%%%%%%%%%%%%%%%%%%%%%%%%
%    Contact details     %
%%%%%%%%%%%%%%%%%%%%%%%%%%
% Real phone/email live in contact.tex, which is gitignored so they are not
% published. A clone of the public repo has only contact.example.tex, hence
% the fallback -- without it a fresh clone would fail to build.
% Defines: \contactdetails

\IfFileExists{contact.tex}
  {\input{contact.tex}}
  {%%%%%%%%%%%%%%%%%%%%%%%%%%
%   Placeholder contact details (tracked by git).
%
%   Used only when contact.tex is missing, so that a fresh clone of the
%   public repo still compiles. To build with real details:
%
%       cp contact.example.tex contact.tex
%
%   then edit contact.tex -- it is gitignored and will not be published.
%%%%%%%%%%%%%%%%%%%%%%%%%%

\newcommand{\contactlinkedin}{\href{https://www.linkedin.com/in/doniyorbek-ibrokhimov}{Linkedin}}

\newcommand{\contactdetails}{%
  \href{https://github.com/doniyorbek-ibrokhimov}{github.com/doniyorbek-ibrokhimov}%
  \contactextra
}
}
, before the document.

% Whether a section is kept on a single page where possible. cv-ios.tex takes
% the default; cv-backend.tex sets this empty, matching its Overleaf original.
\newcommand{\sectionfilbreak}{\filbreak}

% \interests is deliberately NOT defined here. Each variant must define it, so
% a missing definition is a build error rather than a placeholder that could
% slip into a PDF you send out. See cv-ios.tex / cv-backend.tex.

%%%%%%%%%%%%%%%%%%%%%%%%%%
%  Template Definitions  %
%%%%%%%%%%%%%%%%%%%%%%%%%%

\newcommand{\lineunder}{\vspace*{-8pt} \\ \hspace*{-6pt} \hrulefill \\ \vspace*{-15pt}}
\newcommand{\name}[1]{\begin{center}\textsc{\Huge#1}\\\end{center}}
\newcommand{\program}[1]{\begin{center}\textsc{#1}\end{center}}
\newcommand{\contact}[1]{\begin{center}\color{contactgray}{\small#1}\end{center}}

\newenvironment{tabbedsection}[1]{
  \begin{list}{}{
      \setlength{\itemsep}{0pt}
      \setlength{\labelsep}{0pt}
      \setlength{\labelwidth}{0pt}
      \setlength{\leftmargin}{\tabin}
      \setlength{\rightmargin}{\tabin}
      \setlength{\listparindent}{0pt}
      \setlength{\parsep}{0pt}
      \setlength{\parskip}{0pt}
      \setlength{\partopsep}{0pt}
      \setlength{\topsep}{#1}
    }
  \item[]
}{\end{list}}

\newenvironment{nospacetabbing}{
    \begin{tabbing}
}{\end{tabbing}\vspace{-1.2em}}

\newenvironment{resume_header}{}{\vspace{0pt}}


\newenvironment{resume_section}[1]{
  \sectionfilbreak
  \vspace{2\secsep}
  \textsc{\large#1}
  \lineunder
  \begin{tabbedsection}{\secsep}
}{\end{tabbedsection}}

\newenvironment{resume_subsection}[2][]{
  \textbf{#2} \hfill {\footnotesize #1} \hspace{2em}
  \begin{tabbedsection}{0.5\secsep}
}{\end{tabbedsection}}

\newenvironment{subitems}{
  \renewcommand{\labelitemi}{-}
  \begin{itemize}
      \setlength{\labelsep}{1em}
}{\end{itemize}}

\newenvironment{resume_employer}[4]{
  \vspace{\secsep}
  \textbf{#1} \\ 
  \indent {\small #2} \hfill {\footnotesize#3 (#4)}
  \begin{tabbedsection}{0pt}
  \begin{subitems}
}{\end{subitems}\end{tabbedsection}}

%%%%%%%%%%%%%%%%%%%%%%%%%%
%    Contact details     %
%%%%%%%%%%%%%%%%%%%%%%%%%%
% Real phone/email live in contact.tex, which is gitignored so they are not
% published. A clone of the public repo has only contact.example.tex, hence
% the fallback -- without it a fresh clone would fail to build.
% Defines: \contactdetails

\IfFileExists{contact.tex}
  {\input{contact.tex}}
  {%%%%%%%%%%%%%%%%%%%%%%%%%%
%   Placeholder contact details (tracked by git).
%
%   Used only when contact.tex is missing, so that a fresh clone of the
%   public repo still compiles. To build with real details:
%
%       cp contact.example.tex contact.tex
%
%   then edit contact.tex -- it is gitignored and will not be published.
%%%%%%%%%%%%%%%%%%%%%%%%%%

\newcommand{\contactlinkedin}{\href{https://www.linkedin.com/in/doniyorbek-ibrokhimov}{Linkedin}}

\newcommand{\contactdetails}{%
  \href{https://github.com/doniyorbek-ibrokhimov}{github.com/doniyorbek-ibrokhimov}%
  \contactextra
}
}
, before the document.

% Whether a section is kept on a single page where possible. cv-ios.tex takes
% the default; cv-backend.tex sets this empty, matching its Overleaf original.
\newcommand{\sectionfilbreak}{\filbreak}

% \interests is deliberately NOT defined here. Each variant must define it, so
% a missing definition is a build error rather than a placeholder that could
% slip into a PDF you send out. See cv-ios.tex / cv-backend.tex.

%%%%%%%%%%%%%%%%%%%%%%%%%%
%  Template Definitions  %
%%%%%%%%%%%%%%%%%%%%%%%%%%

\newcommand{\lineunder}{\vspace*{-8pt} \\ \hspace*{-6pt} \hrulefill \\ \vspace*{-15pt}}
\newcommand{\name}[1]{\begin{center}\textsc{\Huge#1}\\\end{center}}
\newcommand{\program}[1]{\begin{center}\textsc{#1}\end{center}}
\newcommand{\contact}[1]{\begin{center}\color{contactgray}{\small#1}\end{center}}

\newenvironment{tabbedsection}[1]{
  \begin{list}{}{
      \setlength{\itemsep}{0pt}
      \setlength{\labelsep}{0pt}
      \setlength{\labelwidth}{0pt}
      \setlength{\leftmargin}{\tabin}
      \setlength{\rightmargin}{\tabin}
      \setlength{\listparindent}{0pt}
      \setlength{\parsep}{0pt}
      \setlength{\parskip}{0pt}
      \setlength{\partopsep}{0pt}
      \setlength{\topsep}{#1}
    }
  \item[]
}{\end{list}}

\newenvironment{nospacetabbing}{
    \begin{tabbing}
}{\end{tabbing}\vspace{-1.2em}}

\newenvironment{resume_header}{}{\vspace{0pt}}


\newenvironment{resume_section}[1]{
  \sectionfilbreak
  \vspace{2\secsep}
  \textsc{\large#1}
  \lineunder
  \begin{tabbedsection}{\secsep}
}{\end{tabbedsection}}

\newenvironment{resume_subsection}[2][]{
  \textbf{#2} \hfill {\footnotesize #1} \hspace{2em}
  \begin{tabbedsection}{0.5\secsep}
}{\end{tabbedsection}}

\newenvironment{subitems}{
  \renewcommand{\labelitemi}{-}
  \begin{itemize}
      \setlength{\labelsep}{1em}
}{\end{itemize}}

\newenvironment{resume_employer}[4]{
  \vspace{\secsep}
  \textbf{#1} \\ 
  \indent {\small #2} \hfill {\footnotesize#3 (#4)}
  \begin{tabbedsection}{0pt}
  \begin{subitems}
}{\end{subitems}\end{tabbedsection}}

%%%%%%%%%%%%%%%%%%%%%%%%%%
%    Contact details     %
%%%%%%%%%%%%%%%%%%%%%%%%%%
% Real phone/email live in contact.tex, which is gitignored so they are not
% published. A clone of the public repo has only contact.example.tex, hence
% the fallback -- without it a fresh clone would fail to build.
% Defines: \contactdetails

\IfFileExists{contact.tex}
  {\input{contact.tex}}
  {%%%%%%%%%%%%%%%%%%%%%%%%%%
%   Placeholder contact details (tracked by git).
%
%   Used only when contact.tex is missing, so that a fresh clone of the
%   public repo still compiles. To build with real details:
%
%       cp contact.example.tex contact.tex
%
%   then edit contact.tex -- it is gitignored and will not be published.
%%%%%%%%%%%%%%%%%%%%%%%%%%

\newcommand{\contactlinkedin}{\href{https://www.linkedin.com/in/doniyorbek-ibrokhimov}{Linkedin}}

\newcommand{\contactdetails}{%
  \href{https://github.com/doniyorbek-ibrokhimov}{github.com/doniyorbek-ibrokhimov}%
  \contactextra
}
}
